\begin{table}[t]
\centering
\small
\caption{Characteristics of the seven study metropolitan regions across contrasting Indian climatic biomes (2020--2022 panel, Uber H3 Resolution~8).}
\label{tab:city_biomes}
\begin{tabularx}{\textwidth}{@{}l X r r r l@{}}
\toprule
\textbf{City} & \textbf{K{\"o}ppen Biome} & \textbf{H3 Hexes} & \textbf{Area (km$^2$)} & \textbf{Panel Obs.} & \textbf{Centroid Coordinates} \\
\midrule
Bikaner & Hot arid desert (\textit{BWh}) & 959 & 707 & 34,524 & $28.02^\circ\text{N}, 73.32^\circ\text{E}$ \\
Coimbatore & Semi-arid rainshadow (\textit{BSh}) & 995 & 733 & 35,820 & $11.01^\circ\text{N}, 76.96^\circ\text{E}$ \\
Delhi & Subtropical semi-arid (\textit{Cwa}) & 3,398 & 2,504 & 122,328 & $28.61^\circ\text{N}, 77.20^\circ\text{E}$ \\
Hyderabad & Tropical semi-arid (\textit{BSh}) & 1,505 & 1,109 & 54,180 & $17.38^\circ\text{N}, 78.48^\circ\text{E}$ \\
Indore & Tropical savanna plateau (\textit{Aw}) & 731 & 539 & 26,316 & $22.72^\circ\text{N}, 75.86^\circ\text{E}$ \\
Kolkata & Humid deltaic subtropical (\textit{Aw}) & 992 & 731 & 35,712 & $22.57^\circ\text{N}, 88.36^\circ\text{E}$ \\
Mumbai & Tropical coastal monsoon (\textit{Am}) & 707 & 521 & 25,452 & $19.07^\circ\text{N}, 72.88^\circ\text{E}$ \\
\midrule
\textbf{Total} & \textbf{7 Contrasting Biomes} & \textbf{9,287} & \textbf{6,844} & \textbf{334,332} & \textbf{Pan-India Cohort} \\
\bottomrule
\end{tabularx}
\end{table}
